\subsubsection{Adiabatic Viscous\slash No-Slip Wall}
\label{SEC:BC:Wall:Adiabatic}

The adiabatic viscous\slash no-slip wall boundary condition enforces a
prescribed wall velocity at the boundary surface with zero heat flux through the
boundary, \thatis $\partial T/\partial \normal = 0$.
%
Note, this boundary condition defaults to a stationary wall with zero wall
velocity, but this can be overridden by specifying non-zero values for any of
the boundary condition velocity components.
%
The thermodynamic properties for this boundary condition are just copied from
the interior:
%
\begin{alignat*}{1}
  \rho_\subR & = \rho_\subL \\
  T_\subR & = T_\subL
\end{alignat*}

%%%%%%%%%%%%%%%%%%%%%%%%%%%%%%%%%%%%%%%%%%%%%%%%%%%%%%%%%%%%%%%%%%%%%%%%%%%%%%%%

For most use cases, the prescribed wall velocity is simply zero, \forexample a
stationary flat plate.
%
However, there are certain situations where prescribing a non-zero wall velocity
is needed; a canonical example would be the classic lid-driven cavity problem,
and a more real-world example would be a spinning hub in a turbomachinery
problem.
%
If the prescribed wall velocity is not rotational, the following components of
\BCinptext{} directly set the wall velocity for the boundary condition:
%
\begin{InputExample}
  \item[{\BCinpBC[nmlBC::BCInput/VelocityX]{VelocityX}}]
  \item[{\BCinpBC[nmlBC::BCInput/VelocityY]{VelocityY}}]
  \item[{\BCinpBC[nmlBC::BCInput/VelocityZ]{VelocityZ}}]
\end{InputExample}
%
If the prescribed wall velocity is rotational, the input array
%
\begin{InputExample}
\item[{\BCinpBC[nmlBC::BCInput/RotationRate]{Rotation\_Rate(1:3)}[Rotation\_Rate]}]
\end{InputExample}
%
is used to specify the rotation rate, $\vec{\Omega}$, (in RPMs) for all wall
points in the boundary condition.
%
%The prescribed wall velocity at each wall point depends on the radius
%of the wall point to the axis of rotation, $\vec{r}$, and is given by
%(assuming implicit unit conversions):
%
The prescribed wall velocity at each wall point depends on the radius,
$\vec{r}$, from the axis of rotation to the wall point.
%
Since the axis of rotation for rotating boundaries originates from the
specified 
\inpBC[nmlBC::BCInput/Periodic/RotationCenter]{RotationCenter}%&
      [Periodic\%RotationCenter][],
$\vec{x}_c$, we can also write this in terms of the position vector.
%
The relation is given by (assuming implicit unit conversions):
%
\begin{equation*}
  \vec{u}_\subR = \vec{\Omega} \times \vec{r}
                = \vec{\Omega} \times \left( \vec{x} - \vec{x}_c \right)
\end{equation*}
%
Note, this is different from prescribing a non-rotational wall velocity where
all wall points are set to the same velocity.

%%%%%%%%%%%%%%%%%%%%%%%%%%%%%%%%%%%%%%%%%%%%%%%%%%%%%%%%%%%%%%%%%%%%%%%%%%%%%%%%

If using a moving reference frame, the input
%
\begin{InputExample}
  \item[{\BCinpBC[nmlBC::BCInput/IsSpecifiedInMRF]{Is\_Specified\_In\_MRF}}]
\end{InputExample}
%
can be used to simplify the boundary condition inputs for cases involving
multiple solid wall boundaries with velocities dependent on which reference
frame the boundary is viewed from.
%
An example of this would be a turbomachinery fan containing a spinning hub and
rotors --- with a rotation rate of $\vec{\Omega}$ --- inside a stationary case,
which is how this problem is viewed from the inertial\slash fixed\slash absolute
reference frame.
%
However, when viewed from the rotating reference frame, the hub and rotor are
stationary and the case is now rotating at $(-\vec{\Omega})$.
%
Note, the prescribed wall velocity components are specified as the absolute
velocity components.

%%%%%%%%%%%%%%%%%%%%%%%%%%%%%%%%%%%%%%%%%%%%%%%%%%%%%%%%%%%%%%%%%%%%%%%%%%%%%%%%

To illustrate the use of \BCinptext{Is\_Specified\_In\_MRF}, let us continue
with this turbomachinery example of a spinning hub\slash rotor surrounded by a
stationary case, and our simulation is using a rotating reference frame with
the same rotation rate as the spinning hub\slash rotor.
%
Let \BCinptext[1]{} be the boundary condition inputs for the casing,
and \BCinptext[2]{} be for the hub\slash rotor.
%
Using the default values of zero for the three velocity components, the boundary
conditions for this case can be provided as:
%
\begin{InputExample}
  \item[{\BCinptext[1]{Is\_Specified\_In\_MRF}} = \fls]
  \item[{\BCinptext[2]{Is\_Specified\_In\_MRF}} = \tru]
\end{InputExample}
%
One benefit from this approach is that the boundary condition wall velocities
will automatically adjust if the rotation rate of the rotating reference frame
is changed.

%%%%%%%%%%%%%%%%%%%%%%%%%%%%%%%%%%%%%%%%%%%%%%%%%%%%%%%%%%%%%%%%%%%%%%%%%%%%%%%%

For this boundary condition, the temperature gradient is the only right gradient
state treated differently from the other flow-through boundary conditions
(\eref{EQ:BC:FlowThroughGradient}).
%
The adiabatic condition prescribing zero heat flux through the boundary requires
the temperature gradient to be tangential to the boundary surface.
%
This is exactly the same as the temperature gradient boundary condition for the
slip wall in \sref{SEC:BC:Wall:Slip}, which is
%
\begin{equation*}
  \nabla T_\subR = \nabla T_\subL
                 - \left(\nabla T_\subL\cdot\normal\right)\normal
\end{equation*}

%%%%%%%%%%%%%%%%%%%%%%%%%%%%%%%%%%%%%%%%%%%%%%%%%%%%%%%%%%%%%%%%%%%%%%%%%%%%%%%%

%\textcolor{red}{
%%
%  The viscous wall condition, sometimes referred to as a no-slip wall
%  condition, enforces the prescribed wall tangential velocity along the
%  boundary.
%%
%  The wall normal velocity is set to zero, resulting in no flux through the
%  boundary.
%%
%  For a non-moving viscous wall, all components of the wall velocity are zero. 
%%
%  The adiabatic condition also enforces no heat flux through the boundary, so
%  there will be no wall-normal temperature gradient.
%%
%}

%###############################################################################

